% =========================================================
% Order Extensions
% Source: Mark Dean, Order Theory (Brown University, 2015)
% =========================================================

\subsection{Order Extensions}

% ---------------------------------------------------------
% Symmetric and asymmetric parts
% ---------------------------------------------------------

\begin{definitionbox}{Definition (Symmetric and Asymmetric Parts)}
\begin{definition}[Symmetric and Asymmetric Parts]\label{def:sym-asym}
For any binary relation $R$ on $X$, the \emph{symmetric part} $I$ and
\emph{asymmetric part} $P$ are defined by
\[
  xIy \Longleftrightarrow (xRy \wedge yRx),
  \qquad
  xPy \Longleftrightarrow (xRy \wedge \neg yRx).
\]
Then $R=P\cup I$.
\end{definition}
\end{definitionbox}

\begin{remark*}[Standard quantified statement]
$\forall x,y\in X,\; xIy\leftrightarrow(xRy\wedge yRx)$ and
$\forall x,y\in X,\; xPy\leftrightarrow(xRy\wedge\neg yRx)$.
\end{remark*}

\begin{remark*}[Predicate reading]
$\operatorname{SymmetricPart}(I,R,X)$ records the mutually related pairs of
$R$, and $\operatorname{AsymmetricPart}(P,R,X)$ records the pairs whose reverse
does not belong to $R$.
\end{remark*}

\begin{remark*}[Interpretation]
Every relation decomposes into a part where both directions hold and a part
where only one direction holds.
\end{remark*}

\begin{remark*}[Exposition]
In decision theory, $R$ often encodes a weak preference relation. The symmetric
part is then indifference, and the asymmetric part is strict preference. The
same decomposition is useful for any relation that distinguishes weak and
strict comparison.
\end{remark*}

\begin{remark*}[Exposition]
Preorders and linear orders are defined elsewhere. A preorder is reflexive and
transitive, while a linear order is a total partial order. This section uses
those definitions rather than restating them.
\end{remark*}

\begin{remark*}[Exposition]
A preorder need not be antisymmetric: two distinct elements may satisfy $xRy$
and $yRx$ simultaneously. Adding antisymmetry promotes a preorder to a partial
order. Every partial order is a preorder, but not conversely.
\end{remark*}

\begin{dependencies}
\begin{itemize}
  \item \hyperref[def:relation]{Relation}
  \item \hyperref[def:union]{Union}
  \item \hyperref[def:set-difference]{Set Difference}
  \item \hyperref[def:preorder]{Preorder}
  \item \hyperref[def:total-order]{Linear Order}
\end{itemize}
\end{dependencies}

% ---------------------------------------------------------
% Maximal elements vs. upper bounds
% ---------------------------------------------------------

\begin{definitionbox}{Definition (Maximal Elements and Upper Bounds)}
\begin{definition}[Maximal Elements and Upper Bounds]\label{def:maximal-upper}
Let $\succsim$ be a binary relation on $X$. The set of \emph{maximal elements}
of $X$ is
\[
  \mathrm{Max}(X,\succsim)
  = \{x\in X\mid \text{there is no } y\in X \text{ with } y\neq x
      \text{ and } y\succsim x\}.
\]
The set of \emph{upper bounds} of $X$ is
\[
  M(X,\succsim)
  = \{x\in X\mid \forall y\in X,\; x\succsim y\}.
\]
\end{definition}
\end{definitionbox}

\begin{remark*}[Standard quantified statement]
$x\in\mathrm{Max}(X,\succsim)\leftrightarrow
x\in X\wedge\neg(\exists y\in X,\; y\neq x\wedge y\succsim x)$, and
$x\in M(X,\succsim)\leftrightarrow x\in X\wedge\forall y\in X,\; x\succsim y$.
\end{remark*}

\begin{remark*}[Predicate reading]
$\operatorname{MaximalElement}(x,X,\mathsf{OrderedSet}(X,\succsim))$ means that
nothing in $X$ lies strictly above $x$, while
$\operatorname{UpperBound}(x,X,\mathsf{OrderedSet}(X,\succsim))$ means that
$x$ lies above every element of $X$.
\end{remark*}

\begin{remark*}[Interpretation]
Maximality forbids a strictly better element; being an upper bound requires
being at least as high as every element.
\end{remark*}

\begin{remark*}[Exposition]
A greatest element is always maximal when it exists, but a maximal element need
not be greatest. In a linear order, maximal and greatest coincide because every
pair of elements is comparable.
\end{remark*}

\begin{remark*}[Exposition]
The transitive closure $T(R)$ of a binary relation $R$ on $X$ is the smallest
transitive relation containing $R$, as defined in the Relations chapter.
Equivalently, it records the pairs connected by finite $R$-chains. This
section uses that existing definition without restating it.
\end{remark*}

\begin{dependencies}
\begin{itemize}
  \item \hyperref[def:relation]{Relation}
  \item \hyperref[def:upper-bound]{Upper Bound}
  \item \hyperref[def:strict-order]{Strict Order}
  \item \hyperref[def:relation-transitive-closure]{Transitive Closure}
  \item \hyperref[def:transitive]{Transitive Relation}
\end{itemize}
\end{dependencies}

% ---------------------------------------------------------
% Extensions of preorders
% ---------------------------------------------------------

\begin{definitionbox}{Definition (Extension of a Preorder)}
\begin{definition}[Extension of a Preorder]\label{def:preorder-extension}
Let $\succsim$ be a preorder on $X$. A preorder $\succsim'$ is an
\emph{extension} of $\succsim$ if
\[
  x \succsim y \;\Rightarrow\; x \succsim' y,
  \qquad
  x \succ y \;\Rightarrow\; x \succ' y.
\]
A \emph{complete extension} is an extension that is also complete.
\end{definition}
\end{definitionbox}

\begin{remark*}[Standard quantified statement]
$\operatorname{PreorderExtension}(\succsim',\succsim,X)\leftrightarrow
\operatorname{Preorder}(\succsim',X)\wedge
\forall x,y\in X,\; (x\succsim y\to x\succsim' y)\wedge
(x\succ y\to x\succ' y)$.
\end{remark*}

\begin{remark*}[Predicate reading]
$\operatorname{PreorderExtension}(\succsim',\succsim,X)$ means that the new
preorder keeps every weak comparison and every strict comparison from the old
preorder.
\end{remark*}

\begin{remark*}[Interpretation]
An extension fills in undecided comparisons without reversing the original
strict order information.
\end{remark*}

\begin{remark*}[Exposition]
The purpose of an extension is to promote a partial ranking to a more complete
ranking while preserving the original judgements.
\end{remark*}

\begin{dependencies}
\begin{itemize}
  \item \hyperref[def:preorder]{Preorder}
  \item \hyperref[def:relation]{Relation}
  \item \hyperref[def:sym-asym]{Symmetric and Asymmetric Parts}
\end{itemize}
\end{dependencies}

% ---------------------------------------------------------
% Szpilrajn's Theorem
% ---------------------------------------------------------

\begin{theorembox}{Theorem (Szpilrajn, 1930)}
\begin{theorem}[Szpilrajn Extension Theorem]\label{thm:szpilrajn}
For any nonempty set $X$ and partial order $\succsim$ on $X$, there exists a
linear order that is an extension of $\succsim$.
\hyperref[prf:szpilrajn]{\textit{Go to proof.}}
\end{theorem}
\end{theorembox}

\begin{remark*}[Standard quantified statement]
$\forall X\neq\varnothing,\;\forall \succsim,\;
\operatorname{PartialOrder}(\succsim,X)\to
\exists \succsim^*,\; \operatorname{TotalOrder}(\succsim^*,X)\wedge
\operatorname{PreorderExtension}(\succsim^*,\succsim,X)$.
\end{remark*}

\begin{remark*}[Predicate reading]
$\operatorname{PreorderExtension}(\succsim^*,\succsim,X)$ says the linear order
$\succsim^*$ extends the given partial order $\succsim$.
\end{remark*}

\begin{remark*}[Interpretation]
Every partial order can be completed to a linear order, assuming the standard
choice principles used in the proof.
\end{remark*}

\begin{remark*}[Exposition]
The usual proof orders the set of all partial-order extensions by inclusion
and applies the Hausdorff Maximum Principle. A maximal extension must be
linear; otherwise an incomparable pair can be added and closed transitively,
contradicting maximality.
\end{remark*}

\begin{remark*}[Exposition]
For infinite sets, this theorem depends on choice principles equivalent to the
Axiom of Choice.
\end{remark*}

\begin{dependencies}
\begin{itemize}
  \item \hyperref[def:partial-order]{Partial Order}
  \item \hyperref[def:total-order]{Linear Order}
  \item \hyperref[def:preorder-extension]{Extension of a Preorder}
  \item \hyperref[thm:hausdorff]{Hausdorff Maximum Principle}
\end{itemize}
\end{dependencies}

\begin{corollarybox}{Corollary (Complete preorder extension)}
\begin{corollary}[Complete Preorder Extension]\label{cor:complete-preorder-extension}
For any nonempty set $X$ and preorder $\succsim$ on $X$, there exists a
complete preorder that is an extension of $\succsim$.
\hyperref[prf:complete-preorder-extension]{\textit{Go to proof.}}
\end{corollary}
\end{corollarybox}

\begin{remark*}[Standard quantified statement]
$\forall X\neq\varnothing,\;\forall \succsim,\;
\operatorname{Preorder}(\succsim,X)\to
\exists \succsim',\; \operatorname{CompletePreorderExtension}(\succsim',\succsim,X)$.
\end{remark*}

\begin{remark*}[Predicate reading]
$\operatorname{CompletePreorderExtension}(\succsim',\succsim,X)$ means that
$\succsim'$ is a complete preorder on $X$ and extends $\succsim$.
\end{remark*}

\begin{remark*}[Interpretation]
Every preorder admits a complete preorder extension after quotienting by
indifference and linearly extending the quotient order.
\end{remark*}

\begin{remark*}[Exposition]
The proof passes to the quotient by the symmetric part, applies Szpilrajn's
theorem to the induced partial order on the quotient, and pulls the resulting
linear order back to $X$.
\end{remark*}

\begin{dependencies}
\begin{itemize}
  \item \hyperref[def:preorder]{Preorder}
  \item \hyperref[def:preorder-extension]{Extension of a Preorder}
  \item \hyperref[def:quotient-set]{Quotient Set}
  \item \hyperref[def:sym-asym]{Symmetric and Asymmetric Parts}
  \item \hyperref[thm:szpilrajn]{Szpilrajn Extension Theorem}
\end{itemize}
\end{dependencies}

% ---------------------------------------------------------
% Only Weak Cycles (OWC)
% ---------------------------------------------------------

\begin{definitionbox}{Definition (Only Weak Cycles)}
\begin{definition}[Only Weak Cycles]\label{def:owc}
Let $\succsim$ be a binary relation on $X$ with asymmetric part $\succ$ and
transitive closure $T(\succsim)$. The relation $\succsim$ satisfies
\emph{only weak cycles} (OWC) if
\[
  xT(\succsim)y \;\Rightarrow\; \neg(y \succ x).
\]
\end{definition}
\end{definitionbox}

\begin{remark*}[Standard quantified statement]
$\operatorname{OnlyWeakCycles}(\succsim,X)\leftrightarrow
\forall x,y\in X,\; xT(\succsim)y\to\neg(y\succ x)$.
\end{remark*}

\begin{remark*}[Predicate reading]
$\operatorname{OnlyWeakCycles}(\succsim,X)$ means that no finite chain of weak
comparisons can be reversed by a strict comparison.
\end{remark*}

\begin{remark*}[Interpretation]
OWC rules out strict reversals around cycles generated by weak comparisons.
\end{remark*}

\begin{remark*}[Exposition]
OWC prohibits a chain
$x_1 \succsim x_2 \succsim \cdots \succsim x_n$ together with a strict
reversal $x_n\succ x_1$. Such a configuration cannot survive inside a complete
preorder extension.
\end{remark*}

\begin{dependencies}
\begin{itemize}
  \item \hyperref[def:relation]{Relation}
  \item \hyperref[def:sym-asym]{Symmetric and Asymmetric Parts}
  \item \hyperref[def:relation-transitive-closure]{Transitive Closure}
\end{itemize}
\end{dependencies}

\begin{propositionbox}{Proposition (Only Weak Cycles Characterization)}
\begin{proposition}[Only Weak Cycles Characterization]\label{prop:owc-characterization}
A binary relation $\succsim$ on a nonempty set $X$ can be extended to a
complete preorder if and only if it satisfies only weak cycles.
\hyperref[prf:owc-characterization]{\textit{Go to proof.}}
\end{proposition}
\end{propositionbox}

\begin{remark*}[Standard quantified statement]
$\forall X\neq\varnothing,\;\forall \succsim,\;
(\exists \succsim',\operatorname{CompletePreorderExtension}(\succsim',\succsim,X))
\leftrightarrow \operatorname{OnlyWeakCycles}(\succsim,X)$.
\end{remark*}

\begin{remark*}[Predicate reading]
$\operatorname{OnlyWeakCycles}(\succsim,X)$ is exactly the obstruction-free
condition for extending $\succsim$ to a complete preorder.
\end{remark*}

\begin{remark*}[Interpretation]
Complete preorder extensions exist precisely when weak cycles do not contain
strict reversals.
\end{remark*}

\begin{remark*}[Exposition]
If a strict reversal exists along a weak chain, every extension must preserve
the chain and the strict reversal, producing a contradiction. Conversely, when
only weak cycles hold, the relation can be extended by applying the extension
principle to the preorder generated by its transitive closure.
\end{remark*}

\begin{remark*}[Exposition]
Szpilrajn's theorem relies on Zorn's Lemma or the Hausdorff Maximum Principle,
which are equivalent to the Axiom of Choice. The next section develops these
equivalences directly.
\end{remark*}

\begin{dependencies}
\begin{itemize}
  \item \hyperref[def:owc]{Only Weak Cycles}
  \item \hyperref[def:preorder-extension]{Extension of a Preorder}
  \item \hyperref[cor:complete-preorder-extension]{Complete Preorder Extension}
\end{itemize}
\end{dependencies}
